\documentclass [12pt,a4paper]{article}  %Indicamos el tipo de documento.
\usepackage[spanish]{babel} %Para que lea en Español.
\usepackage{graphicx}
\usepackage[utf8]{inputenc} % Para que acepte acentos.
 %\usepackage{Mathtools} Para que acepte las formulas matemáticas
\begin{document} %Para que enpiece el documento.
\title {Informe sobre $\pi$} %Titulo del documento
\author{Yoselin Armas Ramos} %Autor 
\date{ 11 de Abril de 2014} %Fecha
\maketitle
\begin{abstract}
El objetivo de esta práctica es hacer un informe en latex sobre lo que hemos aprendido del número $\pi$.
Para ello vamos a retomar toda la información dada en prácticas anteriores.
\end{abstract} 
\section{Facultad de Matemáticas} 
El informe debe contener lo siguiente: 
\begin{enumerate}
\item Título
\item Autor
\item Fecha
\item Resumen
\item Secciones y subsecciones
\item Pies de páginas
\item Gráficos y referencias a ellos.
\item Tablas y referencias a ellas.
\item Bibliografía y citas.
\end{enumerate}

\footnote{Primera nota a pie de página }

\section{¿QUÉ ES EL NÚMERO $\pi$?} %Primera sección
El número pi es un número irracional y una de las constantes matemáticas más importantes. 
Se emplea frecuentemente en Matematicas,Física e Ingeniería. El valor aproximado de $\pi$ es 3'14159265358979323846...
Es importante mencionar que es la relación entre la longitud de una circunferencia y su diámetro.
\subsection{Historia del número  $\pi$ }
 Aquí va la historia del número $\pi$  
\subsection{ Aproximación del número $\pi$} % segunda subsección
El número $\pi$ se puede aproximar numéricamente mediante la fórmula siguiente :

$\sum_{i=1}^{n} f(xi)$ con $\f(x)= 4/(1+x^2)$ y $\xi=(i-1/2)/n $ , para i=1,...n.


Además, un mayor número de intervalos hace mayor la precisión en los decimales de la aproximación de $\pi$.
Por lo tanto,se concluye que n ha de constituir un parámetro de la función. 

\section{¿Qué uso le hemos dado al número $\pi$?} %segunda sección

Hemos utilizado el número $\pi$ para realizar varias funciones en Python y aproximaciones durante varias semanas.
A continuación veremos un ejemplo.

\subsection{Ejemplos} 
 Aquí van ejemplos. 
% primera subsección de la sección dos.

\subsection{SUBSECCIÓN FINAL $\pi$} %segunda subsección de la sección dos.
En esta sección vamos a incluir un gráfico.

%incluimos un gráfico de la ULL
\includegraphics{logo.eps}[scale=0.15]

\end{document}
